\chapter{Revisión bibliográfica}
\label{cap:revision}

\section{Pronóstico de capturas pesqueras}

La pregunta de investigación planteada busca determinar el mejor modelo para predecir el volumen de capturas pesqueras en Baja California a partir de variables oceanográficas. Para ello, es pertinente conocer algunos estudios previos donde se hayan aplicado técnicas de modelado y pronóstico similares en contextos pesqueros. A continuación, se presentan trabajos similares que sirvieron de guía y referencia:


Yoon et al. \cite{yoon2020} compararon modelos de máquina de vectores de soporte, bosques aleatorios y redes neuronales profundas para predecir capturas de Caballa, Anchoa y Calamar en zonas de pesca de Corea del Sur, utilizando datos de satélite, oceánicos y climáticos integrados en una cuadrícula geográfica con tamaño unitario de 15kmx15km. Contaban con 4,000 registros de captura entre 2014 y 2016. Los mejores resultados los obtuvieron con DNN, sugiriendo su utilidad para tareas predictivas pesqueras con datos limitados.


Otros autores como Zhao et al.\cite{zhao2021} y Adibi et al.\cite{adibi2020} han explorado técnicas como redes neuronales convolucionales y bosques aleatorios en contextos pesqueros. El trabajo de Zhao et al. se desarrolló en el Mar de China Oriental; utilizaron datos de un Sistema de Monitoreo de Buques para estudiar su comportamiento a través de las trayectorias que siguen al pescar. Su análisis espaciotemporal de las distribuciones del esfuerzo pesquero les ayuda a hacer predicciones a corto plazo; se basan en el descubrimiento de relaciones cronológicas de pesca entre buques para predecir distribuciones, aprovechando las relaciones cronológicas entre buques —el buque que sale primero da la pauta a los buques que salen después—. Usan los comportamientos de las primeras embarcaciones del día como indicadores para las predicciones a corto plazo. A pesar de ser una idea interesante, nosotros no tenemos datos de trayectorias de las embarcaciones, pero podríamos explorar esta posibilidad en el futuro. 

Por otra parte, el trabajo de Adibi et al. fue en el norte del Adriático. Se utilizó una cuadrícula con unidades de 5kmx5km y los datos de desembarque procedían de la lonja de Chioggia. Tomaron en cuenta una distribución ponderada de las capturas, ya que las zonas con más buques pesqueros durante determinados periodos tenían más probabilidades de obtener mayores índices de capturas. Para predecir las capturas por unidad de esfuerzo (CPUE), emplearon el modelo Random Forest, con una previsión ingenua de referencia a modo de comparación. Los resultados mostraron una mejora del 13\% con respecto al modelo de referencia.


Existen también propuestas recientes enfocadas en optimizar arquitecturas de redes recurrentes como LSTM \citep{greff2017,lindemann2021} y de redes de atención \citep{vaswani2023} para procesamiento de series de tiempo, las cuales podrían ser opciones viables dada la naturaleza secuencial de los datos pesqueros; el ensamble de una red recurrente con un modelo de \emph{gradient boosting}, como el que se evalúa en el Capítulo~\ref{cap:resultados}, es una combinación usada en otros dominios de series cortas y ruidosas \citep{shi2022}.


Finalmente, establecer modelos estadísticos autorregresivos como ARIMA y SARIMA como baseline o referencia, ha mostrado ser una estrategia útil en estudios previos como el de Makwinja et al.\cite{makwinja2021} para evaluar el valor predictivo agregado de técnicas más complejas de machine learning y deep learning. El trabajo de Makwinja et al. en el lago Malombe, en Malawi, usa solamente datos de arribos pesqueros, ya que ARIMA es un modelo autorregresivo y para su ajuste no necesita las variables oceanográficas y meteorológicas.


El mecanismo ambiental que resultó decisivo para interpretar estos resultados ---las olas de calor marinas--- se revisa por separado en la Sección~\ref{sec:revision_mhw}.

En conjunto, estos trabajos previos proveen contexto y sustento metodológico para el estudio que se propone. Se evaluarán diversas técnicas de modelado predictivo y se compararán contra métodos estadísticos simples que sirvan de referencia \citep{ho1999,hyndman2021}.


\section{La pesquería de langosta roja en Baja California}
\label{sec:revision_pesqueria}

La langosta roja (\emph{Panulirus interruptus}) se distribuye a lo largo de la costa del
Pacífico de la península de Baja California y sostiene una de las pesquerías ribereñas de mayor
valor del noroeste de México, organizada en cooperativas con zonas de captura asignadas. Su
aprovechamiento está regulado por la Norma Oficial Mexicana NOM-006-SAG/PESC-2016 \citep{nom006pesc2016}, que fija
tallas mínimas, artes de pesca permitidas y la temporada de captura; el diagnóstico del estado
del recurso se publica en la Carta Nacional Pesquera \citep{cnp2025}. La ventana de 15 de
septiembre a 15 de febrero que estructura todo este trabajo proviene de ese marco regulatorio,
no de una decisión de modelado.

Para el pronóstico importa que la especie es sensible a la variabilidad ambiental.
\citet{vega2003} documenta, con una serie de once años frente a la costa central de Baja
California, que las estrategias reproductivas de \emph{P. interruptus} ---medidas por la
proporción de hembras con espermatóforo y de hembras ovígeras--- responden a la temperatura y a
la intensidad de las surgencias. Esa dependencia térmica es el mecanismo que hace plausible
usar variables oceanográficas rezagadas como predictoras de la captura meses después.

En el contexto mexicano existe además un antecedente directo de pronóstico sobre esta misma
pesquería: \citet{hernandez2022} comparan redes neuronales artificiales (NAR y NARX) contra un
modelo ARIMAX para pronosticar el \emph{precio} de la langosta roja mexicana. La estructura del
problema es la misma que aquí ---serie corta, comparación contra un referente autorregresivo---
aunque la variable objetivo es distinta: el precio de mercado en lugar del volumen desembarcado.

\section{Olas de calor marinas y su efecto sobre las pesquerías}
\label{sec:revision_mhw}

Una ola de calor marina (\emph{marine heatwave}, MHW) es un periodo prolongado de temperatura
superficial del mar anómalamente alta. \citet{hobday2016} propusieron la definición que se ha
vuelto estándar ---al menos cinco días consecutivos por encima del percentil 90 de una
climatología de referencia--- y \citet{hobday2018} la extendieron con un esquema de categorías
(moderada, fuerte, severa, extrema) que permite comparar eventos entre regiones. Esa definición
es la que se implementa en el Capítulo~\ref{cap:datos}.

Estos eventos se han vuelto más frecuentes y más largos: \citet{oliver2018} documentan un
aumento sostenido durante el último siglo, y \citet{frolicher2018} proyectan su intensificación
bajo calentamiento global. El Pacífico nororiental ---la región de este trabajo--- vivió uno de
los episodios mejor documentados, el llamado ``Blob'' de 2014-2016, cuya persistencia
interanual analizan \citet{dilorenzo2016} y cuyos efectos biológicos, desde el plancton hasta
los mamíferos marinos, sintetizan \citet{cavole2016}.

El vínculo con la captura está documentado a dos escalas. Globalmente, \citet{free2019}
estiman que el calentamiento histórico ya redujo la productividad pesquera máxima sostenible, y
\citet{smith2021} revisan los impactos socioeconómicos de las MHW sobre las comunidades que
dependen de la pesca. Localmente ---y es el hallazgo que reorientó este trabajo---
\citet{villasenor2024} documentan que las olas de calor marinas redujeron las capturas de
langosta, erizo y pepino de mar en la península de Baja California entre un 15\% y un 58\%
durante las últimas dos décadas, con impactos mayores cerca de los quiebres biogeográficos como
San Quintín. Esa evidencia es la que justifica incorporar un índice de MHW como covariable y la
que ofrece una hipótesis concreta sobre la caída de la temporada 2021-2022.

\section{Modelos globales y agrupamiento parcial}
\label{sec:revision_globales}

Cuando cada serie individual es corta, la alternativa a entrenar un modelo por serie es
entrenar \emph{uno solo} sobre todas ellas ---un \emph{modelo global}--- de modo que los
parámetros se estimen con la unión de las observaciones y las series cortas se apoyen en las
largas. \citet{montero2021} muestran que los modelos globales no son mejores por ser más
complejos sino porque el conjunto agrupado permite ajustar modelos con más capacidad sin
sobreajustar, y que su ventaja no requiere que las series sean ``parecidas''.
\citet{petropoulos2022} documentan el mismo fenómeno en la práctica del pronóstico. Este es
exactamente el régimen de este trabajo: decenas de series de unas pocas temporadas cada una.
La contraparte es que las series de mayor escala dominan la función de pérdida, un problema
que se aborda normalizando la variable objetivo (Sección~\ref{sec:met_global}).

\section{Cuantificación de incertidumbre y predicción conformal}
\label{sec:revision_conformal}

La regresión cuantílica \citep{koenker1978} estima directamente los cuantiles condicionales de
la respuesta en lugar de su media, lo que permite construir intervalos sin suponer una forma
de la distribución del error. Sus intervalos, sin embargo, no tienen garantía de cobertura en
muestra finita: el modelo puede estar mal calibrado.

La \emph{predicción conformal} \citep{vovk2005,angelopoulos2023} corrige eso: a partir de un
conjunto de calibración disjunto del de entrenamiento, ajusta los intervalos para garantizar
cobertura marginal al nivel nominal bajo el único supuesto de intercambiabilidad.
\citet{romano2019} combinan ambas ideas en la \emph{regresión cuantílica conformalizada} (CQR),
que hereda la adaptabilidad de los cuantiles condicionales y añade la garantía de cobertura;
es el método que se adopta en este trabajo. Cuando los cuantiles estimados se cruzan ---un
artefacto frecuente al estimarlos por separado--- se aplica el reordenamiento monótono de
\citet{chernozhukov2010}. Para evaluar pronósticos probabilísticos se usa el CRPS
\citep{gneiting2007}, una regla de puntuación propia que premia simultáneamente calibración y
nitidez.

\section{Interpretabilidad}
\label{sec:revision_shap}

Con decenas de covariables derivadas y pocas temporadas, saber \emph{cuáles} variables usa el
modelo es tan importante como su error. Los valores SHAP \citep{lundberg2017} atribuyen la
predicción a cada variable con una formulación basada en la teoría de juegos cooperativos, y
permiten podar el conjunto de variables con un criterio explícito en lugar de por prueba y
error.
